\section*{Intellectual Foundations, Standards, and Further Reading}

The accounting standards, regulatory guidance, and academic literature below establish the conceptual and legal frameworks governing financial reporting, revenue recognition, and unit economics.

\paragraph{Accounting frameworks and regulatory guidance.}
\begin{itemize}
  \item IFRS Foundation, \emph{Conceptual Framework for Financial Reporting}, revised 2018: definitions of assets, liabilities, equity, income, and expenses; objective of financial reporting.
  \item U.S. Securities and Exchange Commission, \emph{Beginners' Guide to Financial Statements}: accessible overview of the financial statements and the operational distinction between income and cash.
  \item IFRS Foundation, \emph{IFRS 15: Revenue from Contracts with Customers}: core 5-step model for identifying performance obligations and recognizing revenue as goods or services transfer.
  \item Government of India, Ministry of Corporate Affairs, \emph{Indian Accounting Standard (Ind AS) 115: Revenue from Contracts with Customers}: applicable statutory revenue principles in India.
  \item IFRS Foundation, \emph{IAS 7 Statement of Cash Flows}: mandatory classification of operating, investing, and financing cash flows.
  \item IFRS Foundation, \emph{IAS 1 Presentation of Financial Statements}: presentation standards and accrual accounting requirements.
  \item OpenStax, \emph{Principles of Accounting, Volume 1: Financial Accounting}: fundamental instructional reference on debit and credit entries and statement interrelationships.
\end{itemize}

\paragraph{Customer economics and business-model literature.}
\begin{itemize}
  \item Peter S. Fader, Bruce G. S. Hardie, and Ka Lok Lee (2005), `` `Counting Your Customers' the Easy Way: An Alternative to the Pareto/NBD Model,'' \emph{Marketing Science}, 24(2), 275 to 284.
  \item Alexander Osterwalder and Yves Pigneur (2010), \emph{Business Model Generation: A Handbook for Visionaries, Game Changers, and Challengers}, John Wiley \& Sons.
  \item Michael E. Porter (1985), \emph{Competitive Advantage: Creating and Sustaining Superior Performance}, Free Press, on value chains, cost drivers, and operational linkages.
  \item David J. Teece (2010), ``Business Models, Business Strategy and Innovation,'' \emph{Long Range Planning}, 43(2 to 3), 172 to 194.
  \item Founder-supplied companion essay: \emph{The Theory of Pricing and Value Capture} (October 2026), addressing pricing structures, contribution margins, customer acquisition economics, founder labor normalization, and cash discipline.
\end{itemize}

\paragraph{Interpretation warning.}
Managerial contribution, CAC, LTV, normalized founder labor, and unit economics depend on internal operational definitions. They are not interchangeable with statutory net income, recognized revenue, or audited accounting measures under GAAP or IFRS. Where any example uses a straight-line revenue or depreciation assumption, that assumption is made solely to explain the underlying mechanism; real contracts must be analyzed individually under applicable entity laws and tax regulations.
